\begin{table}[h!]
\centering
\scriptsize
\caption{Approximate power of the observed BOSQUE phototype-group sizes for
binary performance contrasts.}
\label{tab:app-orp-current-power}
\setlength{\tabcolsep}{3.2pt}
\renewcommand{\arraystretch}{1.10}
\begin{tabular}{@{}lrrrrrr@{}}
\toprule
\textbf{Eligible units} & $n_{\mathrm{light}}$ & $n_{\mathrm{dark}}$ &
\textbf{Power, $\delta=0.10$} & \textbf{Power, $\delta=0.15$} &
\textbf{MDD, 80\%} & \textbf{MDD, 90\%} \\
\midrule
All observations (accuracy-like) & 105 & 46 & 29.6\% & 53.9\% & 0.212 & 0.250 \\
Malignant lesions (recall) & 77 & 21 & 17.6\% & 31.5\% & 0.307 & 0.361 \\
Benign lesions (specificity) & 28 & 25 & 15.0\% & 26.2\% & 0.348 & 0.408 \\
\bottomrule
\end{tabular}

\vspace{2pt}
\parbox{0.98\textwidth}{\footnotesize
\textit{Notes:} MDD denotes the minimum detectable absolute light--dark gap.
Calculations assume a light-condition performance of $0.85$, a dark-condition
performance of $0.85-\delta$, a two-sided $\alpha=0.05$ test, independent
eligible lesions, and the large-sample two-proportion approximation. The table
is a prospective planning diagnostic; it is not a reanalysis of the bootstrap
contrasts and is not applicable directly to F1-score or AUC--PR.}
\end{table}
